% =====================================================================
%  MODELLO: GLOSSARIO
%
%  Il Glossario è uno solo e sta nella milestone in corso. Si crea
%  copiando TUTTA questa cartella (glossario.tex + capitoli/):
%       RTB/documenti/esterni/glossario/glossario.tex
%       RTB/documenti/esterni/glossario/capitoli/a.tex ...
%  Quando inizia la milestone successiva si SPOSTA (non si copia),
%  insieme agli altri documenti che continuano: GUIDA.md, ricetta 13.
%
%  Le voci NON si scrivono qui ma nei file delle lettere, in capitoli/:
%  una lettera per file (a.tex, b.tex, ...), voci in ordine alfabetico.
%  Per una lettera nuova basta creare il suo file: questo file la
%  include da solo. Guida completa: GUIDA.md (ricetta 16).
% =====================================================================
\documentclass[documento]{gpt}

\titolo{Glossario}
\uso{esterno}
\responsabile{spagnoli}            % chi approva il documento: il Responsabile

% Facoltativo, solo dopo l'assegnazione del capitolato:
% \capitolato{Cn -- Nome del capitolato (Azienda)}

% ---------------------------------------------------------------------
%  REGISTRO DELLE MODIFICHE: qui non si scrive nulla.
%  Le righe le aggiunge il bot dopo il merge della pull request, nel file
%  glossario.registro.tex accanto a questo, in base all'etichetta della PR:
%    contenuto     prima stesura e parti nuove   (0.1.0, 0.2.0, ...)
%    correzione    correzioni piccole             (0.2.1, ...)
%    approvazione  versione approvata             (1.0.0, ...)
%  Autore = chi ha fatto la PR, verificatore = chi l'ha approvata.
%  Approvazione (PR separata, etichetta messa dal Responsabile): nella riga
%  solo il Responsabile come autore, nessun verificatore.
%  stesura-e-approvazione (la mette il Responsabile): le due righe insieme,
%  contenuto (0.1.0) e approvazione (1.0.0), da una sola PR.
%  Regole: GUIDA.md, sezione 5.
% ---------------------------------------------------------------------

\begin{document}

\section*{Introduzione}
Questo documento raccoglie i termini tecnici usati nella documentazione
del progetto. Nei documenti, i termini spiegati qui sono contrassegnati
dalla lettera G a pedice.

% Una lettera per file: se il file non esiste, la lettera si salta.
% Non serve toccare queste righe quando aggiungete una lettera.
\InputIfFileExists{capitoli/a.tex}{}{}
\InputIfFileExists{capitoli/b.tex}{}{}
\InputIfFileExists{capitoli/c.tex}{}{}
\InputIfFileExists{capitoli/d.tex}{}{}
\InputIfFileExists{capitoli/e.tex}{}{}
\InputIfFileExists{capitoli/f.tex}{}{}
\InputIfFileExists{capitoli/g.tex}{}{}
\InputIfFileExists{capitoli/h.tex}{}{}
\InputIfFileExists{capitoli/i.tex}{}{}
\InputIfFileExists{capitoli/j.tex}{}{}
\InputIfFileExists{capitoli/k.tex}{}{}
\InputIfFileExists{capitoli/l.tex}{}{}
\InputIfFileExists{capitoli/m.tex}{}{}
\InputIfFileExists{capitoli/n.tex}{}{}
\InputIfFileExists{capitoli/o.tex}{}{}
\InputIfFileExists{capitoli/p.tex}{}{}
\InputIfFileExists{capitoli/q.tex}{}{}
\InputIfFileExists{capitoli/r.tex}{}{}
\InputIfFileExists{capitoli/s.tex}{}{}
\InputIfFileExists{capitoli/t.tex}{}{}
\InputIfFileExists{capitoli/u.tex}{}{}
\InputIfFileExists{capitoli/v.tex}{}{}
\InputIfFileExists{capitoli/w.tex}{}{}
\InputIfFileExists{capitoli/x.tex}{}{}
\InputIfFileExists{capitoli/y.tex}{}{}
\InputIfFileExists{capitoli/z.tex}{}{}

\end{document}
